\chapter*{Start here: from zero}
\addcontentsline{toc}{chapter}{Start here: from zero}
\markboth{Start here}{Start here}

You do not need to have seen a proof before. This book assumes you can add,
multiply, and tell an even number from an odd one, and it builds everything else.
This short chapter explains \emph{why} a mathematician is not satisfied by
checking examples, lists the few school facts we use, and shows how the symbols
in Chapter~1 read as ordinary English.

\section*{Why prove anything?}

Try this. Compute $n^2+n+41$ for $n=0,1,2,\dots$:
\[
n=0:\ 41,\qquad n=1:\ 43,\qquad n=2:\ 47,\qquad n=3:\ 53,\qquad n=4:\ 61,\ \dots
\]
Every value is a prime number. It stays prime for $n=0,1,\dots,39$, forty cases in
a row. Surely the formula always gives a prime? It does not: at $n=40$,
\[
40^2+40+41 = 40(40+1)+41 = 40\cdot 41 + 41 = 41\cdot 41 = 1681,
\]
which is $41^2$, not prime. Forty correct examples were not a proof, and the
forty-first case broke it. A \emph{proof} is an argument that covers \emph{every}
case at once, so that no later case can surprise you. Contest mathematics and
computer science run on that guarantee: an algorithm proved correct is correct on
inputs you never tested.

\section*{School facts we use}

\begin{description}
\item[Integers.] The whole numbers $\dots,-2,-1,0,1,2,\dots$; the set of all of
      them is written $\mathbb Z$. The positive ones, $1,2,3,\dots$, are the
      \emph{natural numbers}, $\mathbb N$.
\item[Even and odd.] An integer is even if it can be written $2k$ for an
      integer $k$, and odd if it can be written $2k+1$. Every integer is exactly
      one of the two.
\item[Divides.] $a\mid b$ means $b=ak$ for some integer $k$: $3\mid 12$ because
      $12=3\cdot4$; $5\nmid 12$.
\item[Prime.] An integer greater than 1 whose only positive divisors are $1$ and
      itself: $2,3,5,7,11,\dots$. The number $1$ is not prime.
\item[Fractions and real numbers.] $\mathbb Q$ is the set of fractions $a/b$ with
      $a,b$ integers and $b\ne0$; $\mathbb R$ is the set of all numbers on the
      number line.
\end{description}

\section*{Reading the symbols}

Every symbol below is shorthand for an English phrase. Chapter~1 introduces them
one by one; this table is for reference.

\begin{center}
\begin{tabular}{lP{9cm}}
\toprule
\thead{Symbol} & \thead{Read it as}\\
\midrule
$x\in S$ & ``$x$ is an element of $S$'' ($3\in\{1,2,3\}$)\\
$\forall x$ & ``for every $x$''\\
$\exists x$ & ``there is at least one $x$ such that''\\
$P\Rightarrow Q$ & ``if $P$ then $Q$''\\
$P\Leftrightarrow Q$ & ``$P$ if and only if $Q$'': each implies the other\\
$\neg P$ & ``not $P$''\\
$P\wedge Q$, $P\vee Q$ & ``$P$ and $Q$'', ``$P$ or $Q$ (or both)''\\
\bottomrule
\end{tabular}
\end{center}

\section*{How the books of this series fit together}

This is Book~2 of eight. It is the proof-and-counting spine for the others, and it
is a good first book if you plan to continue.

\begin{center}
\begin{tabular}{P{2.4cm}P{9.0cm}}
\toprule
\thead{Book} & \thead{Connection}\\
\midrule
1 Competitive programming & Implements the number theory (Chapters~5--7 here) and counting (Chapters~8--9) as fast code, and uses induction and exchange arguments to justify greedy algorithms.\\
3 Linear algebra & Uses the proof habits of Chapters~2--4: every theorem there is stated, then proved, in the same layout.\\
6 Quant research & Needs the finite probability of Chapter~12.\\
\bottomrule
\end{tabular}
\end{center}

\section*{How to use the boxes in each chapter}

Each chapter begins with \emph{After this chapter you can}, naming what you will
be able to \emph{do}. From Chapter~2, a \emph{Before you start} box lists what
earlier chapters (or other books) the chapter relies on, with a one-sentence
recap, so that you can tell whether to reread something first.

Work the exercises with a pencil. A proof you only read is a proof you have
agreed with; a proof you wrote is one you can use.
